\documentclass[11pt]{article}
\usepackage{fontspec}
\usepackage[french]{babel}
\usepackage{amsmath}
\usepackage{amssymb}
\usepackage{amsthm}
\usepackage{mathrsfs}
\usepackage{tikz-cd}
\usepackage{enumitem}

\frenchsetup{AutoSpacePunctuation=false}

\newcommand{\printedpage}[1]{\marginpar{\footnotesize p.~#1}}
\newcommand{\msp}[1]{\marginpar{\footnotesize [#1]}}
\newcommand{\Ho}{\mathrm{Ho}}
\newcommand{\Cof}{\mathrm{Cof}}

\begin{document}

\printedpage{170}
% head: A. Grothendieck : Dérivateurs, Ch. XII, Édition des Universités de Montpellier II et Paris VII

Tous ces résultats, et la notion de carré $W$-cartésien, se dualisent (moyennant l'axiome idoine de factorisation à droite, cf. page 240) en une théorie des sommes amalgamées $W$-homotopiques, et des carrés $W$-cocartésiens. Cette fois on travaille avec la catégorie
\[
\Phi \;=\; \left(
\begin{tikzcd}[row sep=small, column sep=small]
 & 1 \arrow[dl] \\
0 & \\
 & 2 \arrow[ul]
\end{tikzcd}
\right)
\qquad\quad \overset{j}{\hookrightarrow} \qquad\quad
\tilde{\Phi} \;=\; \left(
\begin{tikzcd}[row sep=small, column sep=small]
 & 1 \arrow[dl] & \\
0 & & 3 \arrow[ul] \arrow[dl] \\
 & 2 \arrow[ul] &
\end{tikzcd}
\right),
\]
$j$ étant une immersion fermée. On a un foncteur restriction
\[
j^* : \Ho_W(\tilde{\Phi}) \longrightarrow \Ho_W(\Phi),
\]
dont on construit un adjoint à gauche
\[
j_! : \Ho_W(\Phi) \longrightarrow \Ho_W(\tilde{\Phi}).
\]
Et on trouve une notion naturelle de
\msp{page 262}
carré $W$-cocartésien, qui paraît irréprochable à tout points de vue, et d'une simplicité parfaite. % [sic] tout points -> tous points, agreement of "tout"

Cela amène à revoir la notion en question dans Cat, où elle avait l'air très astucieuse, profonde et un peu obscure, avec des intégrales et tout ça. C'est le moment de se rendre compte que `ce n'est que ça'. Donc de prouver la

\textbf{Proposition 20.} \textit{Supposons} $\mathcal{M} = \mathrm{Cat}$\textit{, $W$ un localiseur fondamental satisfaisant W(7 bis). Considérons un carré commutatif dans} Cat
\[
Q = \left(
\begin{tikzcd}
S \arrow[r] \arrow[d] & X_1 \arrow[d] \\
X_2 \arrow[r] & X_3
\end{tikzcd}
\right).
\]
\textit{Pour qu'il soit $W$-cartésien} \texttt{[plutôt} $W$\texttt{-cocartésien]} \textit{au sens de la théorie développée ici (pour des catégories de modèles générales), il faut et il suffit que}
\[
\varphi_Q : \quad \mathcal{X}_Q = \left(
\begin{tikzcd}[row sep=small, column sep=small]
 & X_1 \arrow[dl] \\
\int S & \\
 & X_2 \arrow[ul]
\end{tikzcd}
\right) \qquad\quad \longrightarrow \qquad\quad X_3
\]

\texttt{[plutôt}
\[
\varphi_Q : \quad \mathcal{X}_Q = \left(
\begin{tikzcd}[row sep=small, column sep=small]
 & X_1 \\
\int S \arrow[ur] \arrow[dr] & \\
 & X_2
\end{tikzcd}
\right) \qquad\quad \longrightarrow \qquad\quad X_3\text{\texttt{]}}
\]
\printedpage{171}
% head: A. Grothendieck : Dérivateurs, Ch. XII, Édition des Universités de Montpellier II et Paris VII
\textit{soit dans $W$.} (NB L'intérêt de ce critère est qu'il est `intrinsèque' à $Q$, ne faisant pas appel à une factorisation de $\beta_1$ ou de $\beta_2$ \texttt{[où} $\beta_i : S \longrightarrow X_i$\texttt{]} en $S \xrightarrow{\;\in \Cof_W\;} \tilde{X}_i \xrightarrow{\in W} X_i$. Cf. application intéressante dans section V 3.4., p. 146 ff., th. 4 et corollaires.)

\msp{page 263}
\textsc{Démonstration.} Considérons une factorisation de $S \longrightarrow X_1$ en
\[
S \xrightarrow{i_1} \tilde{X}_1 \xrightarrow{p_1} X, % [sic] X -> X_1, the factorisation is of S -> X_1
\]
avec $i_1$ immersion ouverte de Dwyer, $p_1 \in W$. Donc on a un carré commutatif
\[
\tilde{Q} = \left(
\begin{tikzcd}
S \arrow[r, hook, "\mathrm{Dw}"] \arrow[d] & \tilde{X}_1 \arrow[d] \\
X_2 \arrow[r] & X_3
\end{tikzcd}
\right)
\]
et un homomorphisme $\tilde{Q} \xrightarrow{\alpha} Q$ qui est l'identité sur les trois sommets distincts de $\tilde{X}_1$, et qui est $p_1$ en $\tilde{X}_1$. On en déduit un diagramme commutatif
\[
\begin{tikzcd}[row sep=small, column sep=small]
\mathcal{X}_Q \arrow[rr, "\mathcal{X}_\alpha = \lambda"] \arrow[dr, "\varphi_Q"'] & & \mathcal{X}_{\tilde{Q}} \arrow[dl, "\varphi_{\tilde{Q}}"] \\
 & X_3, &
\end{tikzcd}
\]
et on a $\lambda \in W$ (car c'est une $\Phi$-flèche de catégories fibrées sur $\Phi$, qui induit une $W$-équivalence sur chacune de trois fibres.) Donc $\varphi_Q \in W \iff \varphi_{\tilde{Q}} \in W$, d'autre part il est clair que $Q$ est $W$-cocartésien si et seulement si $\tilde{Q}$ l'est. Donc on peut supposer que $S \hookrightarrow X_1$ est \texttt{[une]} immersion ouverte (de Dwyer) (\footnote[134]{On utilise ici W(3) (forme forte), i.e. Loc(3 bis) avec les notations de XVI. Mais on peut s'en passer, cf. plus bas.}). Appliquant ceci
\msp{page 264}
en changeant les rôles de $X_1$ et $X_2$, on voit qu'on peut supposer de plus que $S \longrightarrow X_2$ est également \texttt{[une]} immersion ouverte de Dwyer. Considérons alors
\[
\begin{tikzcd}
S \arrow[r, hook] \arrow[d, hook] \arrow[dr, phantom, "Q'"] & X_1 \arrow[d, hook] & \\
X_2 \arrow[r, hook] & X_1 \amalg_S X_2 \arrow[r, "\pi"] & X_3 \\
 & \mathcal{X}_Q \arrow[u, "\varphi_{Q'}"] \arrow[ur, "\varphi_Q"'] &
\end{tikzcd}
\]
\printedpage{172}
% head: A. Grothendieck : Dérivateurs, Ch. XII, Édition des Universités de Montpellier II et Paris VII
On a $\varphi_Q = \pi\varphi_{Q'}$, d'autre part on sait que $\varphi_{Q'} \in W$ (cor. 3, p. 149). Donc
\[
\pi \in W \qquad\quad \iff \qquad\quad \varphi_Q \in W,
\]
où $\pi \in W$ signifie que $Q$ est $W$-cocartésien, q.e.d.

\msp{page 265}

\section*{8\quad Appendice : Retour sur les axiomes W(7), W(7 bis) (pages 75, 76).}

(\footnote[135]{Reflexion provisoire, supplanté par XVI.}). On a vu (prop. 6, p. 78) que W(7 bis) implique W(7). Montrons une réciproque partielle. % [sic] Reflexion -> Réflexion, missing accent; supplanté -> supplantée, agreement (in footnote 135)

\textbf{Proposition.} \textit{Supposons} W(7 bis) (axiome du carré cocartésien) \textit{satisfait, et plaçons nous sous les conditions de l'énoncé de} W(7), \textit{et plus généralement, dans le cas d'un diagramme cubique}
\[
\begin{tikzcd}[row sep=small, column sep=small]
 & X_2 \arrow[rr] \arrow[dd, "f_2" near end] & & X_3 \arrow[dd, "f_3"] \\
X_0 \arrow[ur] \arrow[rr, crossing over] \arrow[dd, "f_0"'] & & X_1 \arrow[ur] & \\
 & Y_2 \arrow[rr] & & Y_3 \\
Y_0 \arrow[ur] \arrow[rr] & & Y_1 \arrow[ur] \arrow[from=uu, crossing over, "f_1"' near end] &
\end{tikzcd}
\]
\textit{où les deux carrés $X$ et $Y$ (face supérieure et inférieure du cube) sont $W$-cocartésiens. Considérons les conditions suivantes sur ce diagramme}
\begin{enumerate}[label={}]
\item[\textit{a)}] \textit{$f_0 \in W$.}
\item[\textit{b)}] \textit{$f_1$ et $f_2 \in W$.}
\item[\textit{c)}] \textit{$f_3 \in W$.}
\end{enumerate}
\textit{Alors on a les implications}
\[
a \text{ \textit{\&} } b \Longrightarrow c, \qquad\quad b \text{ \textit{\&} } c \Longrightarrow a.
\]
\textit{De plus, si $W$ est stable par facteurs directs (p. ex. $W$ fortement saturé), alors on a aussi}
\[
a \text{ \textit{\&} } c \Longrightarrow b,
\]
\textit{donc la conjonction} \texttt{[de]} \textit{deux quelconques parmi les conditions a) b) c) implique la troisième.}

\textsc{Démonstration.} Utilisant la factorisation
\printedpage{173}
% head: A. Grothendieck : Dérivateurs, Ch. XII, Édition des Universités de Montpellier II et Paris VII
\msp{page 266}
fonctorielle d'une flèche quelconque (ici les $f_i$) en $pi$, avec $i$ immersion ouverte de Dwyer et $p \in W$ (p. 229), on est remené au cas où les $f_i$ sont des immersions ouvertes de Dwyer. Mais pour toute flèche $g : S \longrightarrow T$ qui est dans $\Cof_W$, soient $Z(g)$ son mapping-cone % [sic] remené -> ramené, spelling
\[
\begin{tikzcd}[row sep=small]
S \arrow[rr, "g"] \arrow[dr] & & T \arrow[dr, "\alpha(g)"] & \\
 & e \arrow[rr, "\bar{g}"] & & Z(g),
\end{tikzcd}
\]
et
\[
\alpha(g) : T = \mathrm{but}(g) \longrightarrow Z(g)
\]
la flèche canonique (\footnote[136]{C'est un peu vif, car il n'est pas clair que $X_1' \amalg_{X_0} X_2' \longrightarrow X_1 \amalg_{X_0} X_2$ soit \texttt{[dans]} $W$, sans utiliser déjà un des trois ingrédients de W(7) (savoir $f_0, f_1, f_2 \in W \implies f_3 \in W$). Il semble donc que notre argument ne marche que moyennant une hypothèse supplémentaire de cet ordre \dots Mais il y a le \textit{lemme de Brown,} cf. XVI \S3, notamment p. 20.}). Celle-ci est fonctorielle en $g$, de sorte que le diagramme cubique donné se prolonge en un double cube
\[
\begin{tikzcd}[row sep=small, column sep=small]
 & X_2 \arrow[rr] \arrow[dd, "f_2" near end] & & X_3 \arrow[dd, "f_3"] \\
X_0 \arrow[ur] \arrow[rr, crossing over] \arrow[dd, "f_0"'] & & X_1 \arrow[ur] & \\
 & Y_2 \arrow[rr] \arrow[dd, "\alpha_2" near end] & & Y_3 \arrow[dd, "\alpha_3"] \\
Y_0 \arrow[ur] \arrow[rr, crossing over] \arrow[dd, "\alpha_0"'] & & Y_1 \arrow[ur] \arrow[from=uu, crossing over, "f_1"' near end] & \\
 & Z_2 \arrow[rr] & & Z_3 \\
Z_0 \arrow[ur] \arrow[rr] & & Z_1 \arrow[ur] \arrow[from=uu, crossing over, "\alpha_1"' near end] &
\end{tikzcd}
\]
avec
\[
Z_i = Z(f_i),\ \alpha_i = \alpha(f_i) \qquad\quad \text{pour } i \in \{0,1,2,3\}.
\]
Notons que W(7 bis) implique que si $g \in \Cof_W$,
\msp{page 267}
on a
\[
g \in W \qquad\quad \iff \qquad\quad \bar{g} \in W, \qquad\quad \text{i.e. } Z(g)\ W\text{-asphérique},
\]
en particulier
\[
f_i \in W \qquad\quad \iff \qquad\quad Z_i\ W\text{-asphérique}.
\]
Ainsi, l'énoncé à prouver sur le cube initial se ramène à un énoncé sur le carré
\printedpage{174}
% head: A. Grothendieck : Dérivateurs, Ch. XII, Édition des Universités de Montpellier II et Paris VII
\[
\begin{tikzcd}[row sep=small, column sep=small]
 & Z_1 \arrow[dr] & \\
Z_0 \arrow[ur] \arrow[dr] & & Z_3, \\
 & Z_2 \arrow[ur] &
\end{tikzcd}
\]
savoir que les trois conditions
\begin{enumerate}[label={}]
\item[a$'$)] $Z_0$ $W$-asphérique.
\item[b$'$)] $Z_1$ et $Z_2$ $W$-asphériques.
\item[c$'$)] $Z_3$ $W$-asphérique.
\end{enumerate}
donnant lieu aux implications
\[
\mathrm{a}' \text{ \& } \mathrm{b}' \Longrightarrow \mathrm{c}', \qquad\quad \mathrm{b}' \text{ \& } \mathrm{c}' \Longrightarrow \mathrm{a}',
\]
et que si $W$ est stable par facteurs directs, alors
\[
\mathrm{a}' \text{ \& } \mathrm{c}' \Longrightarrow \mathrm{b}'.
\]

C'est à peu de choses près le corollaire de la page 76, complété à la page 77 (corollaire), à cela près qu'on y suppose les flèches des immersions ouvertes, et le carré cartésien (donc aussi cocartésien) dans Cat, i.e. $Z_3 = Z_1 \cup Z_2$, $Z_0 = Z_1 \cap Z_2$. Mais
\msp{page 268}
ce qui importe pour la validité de cet énoncé, c'est (on s'en doute) que $Z_*$ soit $W$-cocartésien. Donc il faut d'abord prouver ceci.

\textbf{Lemme.} \textit{Si le diagramme cubique p. 265 est tel que les faces $X$, $Y$ soient $W$-cocartésiennes, et les $f_i$ ($i \in \{0,1,2,3\}$) des $W$-cofibrations, alors le carré $Z$ des cofibres des $f_i$ est lui aussi $W$-cocartésien.}

Ce lemme résulte de la théorie du `déploiement des diagrammes cubiques', relatif à une catégorie de modèles $(\mathcal{M}, W)$ admettant une théorie des cofibres $W$-homotopiques, qu'il me faudra bien écrire! (Mais l'existence pour $(\mathrm{Cat}, W)$ d'une telle théorie des cofibres $W$-homotopiques, et par suite de celle des déploiements des diagrammes dans $\mathcal{M}$ de type $(\boldsymbol{\Delta}^1)^I$, $I$ ensemble d'indices fini, ne dépend que de W(7 bis) - on n'a pas à faire appel à W(7).)

\msp{page 269}
Ce `deploiement' est un diagramme infini dont toutes les suites infinies, dans les trois directions des axes de coordonnées, sont des `suites de suspension'. Je reproduis la partie utile du déploiement. % [sic] deploiement -> déploiement, missing accent

\end{document}
